% =====================================================================
%  ANEXO E. Informe de estado de la prueba
% =====================================================================
\section{Informe de estado de la prueba}
\label{anx:informes}

\begin{notanormativa}[Correspondencia con el modelo de procesos]
Corresponde a la actividad \textbf{TMC4 «Reportar»} de \normap{2}. Es el producto periódico del
seguimiento definido en el apartado~\ref{sec:seguimiento} y se dirige a los destinatarios del
Cuadro~\ref{tab:comunicacion}.
\end{notanormativa}

\begin{observacion}[No existía informe de estado]
\obsproblema{La plantilla anterior no preveía ningún reporte periódico: entre una entrega del
curso y la siguiente no quedaba constancia del avance, la cobertura, los incidentes ni el
esfuerzo consumido.}
\obscambio{Formato reutilizable con cinco bloques —identificación, avance, resultados de
ejecución, incidentes y riesgos, recursos, decisiones y acciones—, una copia por periodo.}
\obsfuente{OBS §3 y §6 · PDF §1.3 (TMC4)}
\end{observacion}

\begin{instruccion}
Reproduzca este formato una vez por periodo. Un informe de estado útil no se limita a decir
cuánto se avanzó: contrasta lo previsto con lo real, señala lo que está bloqueado y termina con
decisiones y acciones que tienen responsable y fecha. Si un periodo no requiere acciones,
escríbalo explícitamente en lugar de dejar la sección vacía.
\end{instruccion}

\begin{xltabular}{\textwidth}{|E{4.6cm}|Y|}
\caption{Formato de informe de estado de la prueba.}\label{tab:ficha-informe}\\ \hline
\endfirsthead
\multicolumn{2}{@{}l@{}}{\footnotesize\itshape\tablename~\thetable{} (continuación)}\\ \hline
\endhead
\multicolumn{2}{@{}r@{}}{\footnotesize\itshape continúa en la página siguiente}\\
\endfoot
\endlastfoot
  \multicolumn{2}{|l|}{\cellcolor{uvgris}\textbf{Identificación}} \\ \hline
  Informe número / periodo      & \campo{IE-03 — dd/mm/aaaa a dd/mm/aaaa} \\ \hline
  Elaborado por / fecha         & \campo{Nombre — dd/mm/aaaa} \\ \hline
  Destinatarios                 & \campo{Cuadro~\ref{tab:comunicacion}} \\ \hline
  \multicolumn{2}{|l|}{\cellcolor{uvgris}\textbf{Avance}} \\ \hline
  Avance previsto en el periodo & \campo{\% o hitos comprometidos} \\ \hline
  Avance real alcanzado         & \campo{\% o hitos logrados} \\ \hline
  Desviación y su causa         & \campo{Diferencia y explicación} \\ \hline
  \multicolumn{2}{|l|}{\cellcolor{uvgris}\textbf{Resultados de ejecución}} \\ \hline
  Pruebas ejecutadas            & \campo{n} \\ \hline
  Aprobadas                     & \campo{n} \\ \hline
  Fallidas                      & \campo{n} \\ \hline
  Bloqueadas                    & \campo{n y causa del bloqueo} \\ \hline
  Cobertura alcanzada           & \campo{COB-xx = \% (indique la medida, no sólo el número). Mide lo ejercitado; las pruebas fallidas cuentan igualmente} \\ \hline
  \multicolumn{2}{|l|}{\cellcolor{uvgris}\textbf{Incidentes y riesgos}} \\ \hline
  Incidentes nuevos             & \campo{n; claves de los relevantes} \\ \hline
  Incidentes cerrados           & \campo{n} \\ \hline
  Incidentes abiertos críticos  & \campo{Claves y estado} \\ \hline
  Riesgos materializados        & \campo{Claves} \\ \hline
  Riesgos nuevos identificados  & \campo{Descripción y exposición} \\ \hline
  \multicolumn{2}{|l|}{\cellcolor{uvgris}\textbf{Recursos}} \\ \hline
  Esfuerzo consumido / estimado & \campo{h reales de h estimadas} \\ \hline
  Bloqueos y dependencias       & \campo{Qué impide avanzar y de quién depende} \\ \hline
  \multicolumn{2}{|l|}{\cellcolor{uvgris}\textbf{Decisiones y acciones}} \\ \hline
  Decisiones tomadas            & \campo{Qué se decidió y quién} \\ \hline
  Acciones acordadas            & \campo{Acción — responsable — fecha objetivo} \\ \hline
  ¿Requiere replanificación?    & \campo{Sí/No; si sí, según qué regla del Cuadro~\ref{tab:control}} \\ \hline
\end{xltabular}

El Cuadro~\ref{tab:ficha-informe} agrupa la información en cinco bloques —identificación,
avance, resultados, incidentes y riesgos, recursos, decisiones— de modo que el lector pueda
localizar rápidamente lo que le concierne. Cada informe emitido debería quedar archivado según
el Cuadro~\ref{tab:configuracion}.
